\documentclass[11pt]{article}
\usepackage{amsmath,amssymb,booktabs}
\begin{document}

\section*{Step 221: CTMT Recursion Test on BSD}

\subsection*{BSD Diagnostic State}

The BSD cascade map declares the explored route diagnostic-complete but not proof-complete.  The b-52a record decomposes the Bloch--Kato bridge defect as
\[
\Delta_{\rm BSD}^{\rm BK}
=E_{\rm an/period}+E_{\rm ht/reg}+E_{\rm finite}+\sum_p E_p+E_{\rm det}.
\]
The typed external obligation is to construct component maps
\[
\rho_{\rm an/period},\rho_{\rm ht/reg},\rho_{\rm finite},\rho_p,\rho_{\rm det}
\]
into the Bloch--Kato determinant line and prove all component defects vanish under the declared hypotheses.

\subsection*{Literature Audit}

Bloch--Kato formulate the motivic Tamagawa-number special-value target.  Fontaine--Perrin-Riou refine the cohomological formulation.  Burns--Flach develop determinant functors and the equivariant Tamagawa formulation, including relative \(K_0\) targets.  Their 2006 result proves ETNC for Tate motives in specific abelian settings, not for all elliptic curves.

The p-adic component does not close globally by citation alone.  Kato supplies Euler-system/Selmer bounds and one divisibility-type input.  Skinner--Urban prove Iwasawa-Greenberg main-conjecture cases for large classes under explicit hypotheses.  Howard proves one divisibility in a Heegner-point Iwasawa setting.  These results refine the p-adic gate into further theorem/hypothesis gates rather than closing the BSD determinant-line bridge globally.

\subsection*{Recursion Pattern}

The BSD determinant-line lane exhibits the same structural recursion as the RH CTMT lanes:
\[
\text{external gate }X
\quad\leadsto\quad
\text{more refined gate }Y
\quad\leadsto\quad
\text{hypothesis/control/divisibility gates}.
\]
Here:
\[
\text{BK component maps}
\leadsto
\text{ETNC/BK determinant line}
\leadsto
\text{relative }K_0\text{ and determinant functors}
\leadsto
\text{Iwasawa/Euler-system/control rows}.
\]

\subsection*{Verdict}

\[
\boxed{\mathrm{V\_bsd\_CTMT\_recursion\_verified}}.
\]
The RH-derived CTMT recursion pattern is verified on a second track in structural-recursion form.  BSD is determinant-line/component-map terminal rather than literal RH matrix-element terminal, but the cascade behavior is the same: resolving a named external theorem exposes deeper typed gates.

\end{document}
